\section*{Problem 1}

已知矩阵
\[
    A = \begin{bmatrix}
        1 & 1 \\
        1 & -1 \\
        1 & 0
    \end{bmatrix}.
\]
由于矩阵 $A$ 列满秩，其 Moore-Penrose 广义逆 $A^+$ 可以通过公式
\[
    A^+ = \left(A^\top A\right)^{-1}A^\top
\]
来计算。请写出详细的计算过程，求出 $A^+$。

\section*{Problem 2}

假设线性方程组 $A\bm{x} = \bm{b}$ 有解。利用广义逆的通解性质可知，该方程组的解可以统一表示为
\[
    \bm{x} = A^+\bm{b} + \left(I - A^+A\right)\bm{z},
\]
其中 $\bm{z}$ 为任意向量，且 $A^+$ 为 $A$ 的 Moore-Penrose 广义逆。请证明：
\[
    \bm{x}^* = A^+\bm{b}
\]
是该方程组所有解中二范数（模长）最小的解。

\section*{Problem 3}

已知 $A^+$ 是矩阵 $A$ 的 Moore-Penrose 广义逆，满足以下四条定义性质：
\[
    \begin{aligned}
        &\text{(1)}\quad AA^+A = A; &
        &\text{(2)}\quad A^+AA^+ = A^+; \\
        &\text{(3)}\quad \left(AA^+\right)^\top = AA^+; &
        &\text{(4)}\quad \left(A^+A\right)^\top = A^+A.
    \end{aligned}
\]
请利用上述四条性质，证明 $\left(A^+\right)^\top$ 也是 $A^\top$ 的 Moore-Penrose 广义逆，从而得出结论
\[
    \left(A^\top\right)^+ = \left(A^+\right)^\top.
\]
